\section{Introduction}\label{sec:intro}

Public debate, security funding and policy on antisemitism are argued in
counts. In the FBI's hate-crime statistics for 2025, law enforcement
agencies reported 1{,}528 single-bias incidents motivated by anti-Jewish bias
\citep{fbi2026hate,adl2026fbi}. For the same year, the Anti-Defamation
League's annual Audit tabulated 6{,}274 antisemitic incidents of harassment,
vandalism and assault, criminal and non-criminal alike, after a record 9{,}354
in 2024 \citep{adl2026audit,adl2025audit}. Figures like these are routinely
set side by side, but they are not rival estimates of one quantity. Each is a
\emph{list}: an enumeration of the incidents that happened to pass through a
particular set of filters. The number that matters, the total number of
in-scope incidents including those that appear on no list, is never observed.
Criminologists call it the \emph{dark figure} \citep{pezzella2019}. This paper
is about estimating it.

\IfFileExists{figures/fig-funnel.tex}{% =====================================================================
%  Figure: the reporting funnel (LEAD).  Schematic -- widths are NOT data.
%  Requires the house colours defined in preamble.tex.
% =====================================================================
\begin{figure}[t]
\centering
\begin{tikzpicture}[x=1cm,y=1cm,
    stage/.style={font=\footnotesize\bfseries,text=LabNavy,align=center,
                  anchor=south,inner sep=1pt,text width=2.3cm},
    rho/.style={font=\small,text=LabNavy,inner sep=1pt},
    src/.style={font=\scriptsize\itshape,text=LabGrey,align=center,inner sep=1pt,text width=2.3cm},
    leaklab/.style={font=\scriptsize,text=LabGrey,align=center,anchor=north,
                    inner sep=1pt,text width=2.1cm},
    listline/.style={line width=0.9pt,LabNavyMid},
    listlab/.style={font=\scriptsize,text=LabNavy,anchor=west,inner sep=1pt},
  ]
  % ---------------- geometry (schematic) ----------------
  \def\bw{0.25}   % half bar width
  \def\rin{0.22}  % inner radius of the leaks
  \def\yend{-4.55}
  % x positions and heights of the five stages
  \def\xA{0}    \def\hA{4.0}
  \def\xB{2.8}  \def\hB{3.0}
  \def\xC{5.6}  \def\hC{2.0}
  \def\xD{8.4}  \def\hD{1.45}
  \def\xE{11.2} \def\hE{1.12}

  % leak macro: #1 = x of source bar, #2 = h source, #3 = h next
  \newcommand{\leak}[3]{%
    \pgfmathsetmacro{\a}{#1+\bw}%
    \pgfmathsetmacro{\rout}{\rin+#2-#3}%
    \pgfmathsetmacro{\xi}{\a+\rin}%
    \pgfmathsetmacro{\xo}{\a+\rout}%
    \pgfmathsetmacro{\xm}{(\xi+\xo)/2}%
    \fill[LabAccent!16] (\a,-#3) arc[start angle=90,end angle=0,radius=\rout]
        -- (\xo,\yend) -- (\xm,\yend-0.28) -- (\xi,\yend)
        -- (\xi,-#2-\rin) arc[start angle=0,end angle=90,radius=\rin] -- cycle;
    \draw[LabAccent!75,line width=0.5pt] (\a,-#3) arc[start angle=90,end angle=0,radius=\rout]
        -- (\xo,\yend) -- (\xm,\yend-0.28) -- (\xi,\yend)
        -- (\xi,-#2-\rin) arc[start angle=0,end angle=90,radius=\rin];
  }

  % ---------------- flows (retained mass) ----------------
  \foreach \xa/\xb/\hb in {\xA/\xB/\hB,\xB/\xC/\hC,\xC/\xD/\hD,\xD/\xE/\hE}{
    \fill[LabNavy!9] (\xa+\bw,0) rectangle (\xb-\bw,-\hb);
  }
  % ---------------- leaks (lost from the police record) ----------------
  \leak{\xA}{\hA}{\hB}
  \leak{\xB}{\hB}{\hC}
  \leak{\xC}{\hC}{\hD}
  \leak{\xD}{\hD}{\hE}

  % ---------------- stage bars ----------------
  \fill[LabNavy]      (\xA-\bw,0) rectangle (\xA+\bw,-\hA);
  \fill[LabNavy!88]   (\xB-\bw,0) rectangle (\xB+\bw,-\hB);
  \fill[LabNavy!80]   (\xC-\bw,0) rectangle (\xC+\bw,-\hC);
  \fill[LabNavy!72]   (\xD-\bw,0) rectangle (\xD+\bw,-\hD);
  \fill[LabNavy!64]   (\xE-\bw,0) rectangle (\xE+\bw,-\hE);

  % ---------------- stage labels ----------------
  \node[stage] at (\xA,0.12) {Incidents\\[-1pt]\normalfont$N$};
  \node[stage] at (\xB,0.12) {Reaches a\\[-1pt]channel\enspace\normalfont$N_V$};
  \node[stage] at (\xC,0.12) {Reported\\[-1pt]to police};
  \node[stage] at (\xD,0.12) {Classified as\\[-1pt]antisemitic};
  \node[stage] at (\xE,0.12) {Submitted\\[-1pt]to the FBI};

  % ---------------- multiplicative factors & identifying sources ----------
  \foreach \xa/\xb/\r/\s in {
      \xA/\xB/\rvic/{victim survey: told\\[-2pt]police or a recording org.},
      \xB/\xC/\rpol/{NCVS-type survey:\\[-2pt]told the police},
      \xC/\xD/\rcls/{police--NGO\\[-2pt]linked overlap},
      \xD/\xE/\rsub/{agency\\[-2pt]participation data}}{
    \node[rho] at ({(\xa+\xb)/2},-0.32) {$\times\,\r$};
    \node[src,anchor=north] at ({(\xa+\xb)/2},-0.52) {\s};
  }
  \node[src,anchor=north,text width=1.2cm] at (\xB,-\hB-0.08) {MSE};

  % ---------------- the recorded count ----------------
  \def\xc{12.05}  % x of the collector that gathers all K lists for MSE
  \draw[listline] (\xE+\bw,-0.5) -- (\xc,-0.5);
  \fill[LabNavyMid] (\xc,-0.5) circle (1.6pt);
  \node[anchor=west,inner sep=1pt] (nfbi) at (\xc+0.12,-0.5)
      {\color{LabNavy}$\nFBI$};
  \node[anchor=north west,font=\scriptsize,text=LabGrey,inner sep=1pt]
      at (nfbi.south west) {list $k{=}1$};

  % ---------------- leak labels ----------------
  \foreach \x/\dh/\r/\t in {
      \xA/1.0/\rvic/{reached no\\[-2pt]channel},
      \xB/1.0/\rpol/{not reported\\[-2pt]to police},
      \xC/0.55/\rcls/{not recorded\\[-2pt]as bias crime},
      \xD/0.33/\rsub/{recorded but\\[-2pt]not submitted}}{
    \pgfmathsetmacro{\xm}{\x+\bw+\rin+\dh/2}
    \node[font=\footnotesize,text=LabAccent,anchor=north,inner sep=1pt]
        (ll) at (\xm,\yend-0.36) {$1-\r$};
    \node[leaklab] at (ll.south) {\t};
  }

  % ---------------- parallel lists (capture from the whole population) ----
  \foreach \y/\k/\name in {1.55/2/{NGO audit},2.05/3/{community security},
                           2.55/4/{campus / portal}}{
    \draw[listline] (\xB,\y) -- (\xc,\y);
    \fill[LabNavyMid] (\xc,\y) circle (1.6pt);
    \node[listlab] at (\xc+0.12,\y) {$n_{\k+}$\enspace\textcolor{LabGrey}{\name}};
  }
  % sampling ticks: every list records only incidents that reach a channel
  \draw[LabNavyMid,line width=0.5pt] (\xB,0.95) -- (\xB,2.55);
  \foreach \y in {1.55,2.05,2.55}{\fill[LabNavyMid] (\xB,\y) circle (1.4pt);}
  \node[font=\scriptsize\itshape,text=LabGrey,anchor=east,align=right,inner sep=1pt]
      at (\xB-0.15,2.05) {parallel lists\\[-2pt]record incidents\\[-2pt]that reach them};

  % referral dependence: police-recorded incidents flow into the NGO audit
  \draw[-{Stealth[length=4pt]},LabAccent,line width=0.8pt,dash pattern=on 2.5pt off 1.5pt]
      ({(\xC+\xD)/2},0.02) -- ({(\xC+\xD)/2},1.49);
  \node[font=\scriptsize\itshape,text=LabAccent,anchor=west,inner sep=1pt]
      at ({(\xC+\xD)/2+0.05},1.0) {referral};

  % MSE closes the loop: all K lists, including list 1 -> linkage -> N_V
  \draw[LabNavy,line width=0.7pt] (\xc,-0.5) -- (\xc,2.8);
  \draw[-{Stealth[length=5pt]},LabNavy,line width=0.7pt,rounded corners=5pt]
      (\xc,2.8) -- (\xc,3.25) -- (\xB,3.25) -- (\xB,2.72);
  \node[fill=white,font=\scriptsize,text=LabNavy,inner sep=3pt]
      at ({(\xB+\xc)/2},3.25)
      {link the $K$ lists, then multiple-systems estimation:\enspace $\widehat N_V$};
  % a victim survey closes the first stage
  \node[font=\scriptsize,text=LabNavy,anchor=south,align=center,inner sep=1pt]
      at (\xA,1.1) {$N=N_V/\rvic$\\[-1pt]\textcolor{LabGrey}{\itshape(survey)}};
\end{tikzpicture}
\caption{\textbf{The reporting funnel.} Each in-scope antisemitic incident
passes four filters before it appears in the FBI's hate-crime count: it reaches
a recording channel, by a report or by the channel's own detection
($\rvic$); it reaches the police ($\rpol$); the police record an antisemitic
bias motivation ($\rcls$); and the agency submits the record to the FBI
($\rsub$), so that $\E[\nFBI\mid U_d]=N_V\rpol\rcls\rsub$ (\cref{prop:funnel}).
Orange streams are incidents lost from the police record at each stage; grey
italics name the auxiliary data that inform each factor
(\cref{tab:funnel-sources}). The other lists record incidents that reach them
through channels other than the police, except through \emph{referral}
(dashed; \cref{prop:referral}). Linking all $K$ lists, the police list
included, and applying multiple-systems estimation estimates $N_V$. Incidents
that reach no channel are on no list; a victim survey that supplies $\rvic$ is
needed to recover them. Widths are schematic, not estimates.}
\label{fig:funnel}
\end{figure}
}{}%
\paragraph{How incidents go missing}
The police count is the end of a chain of conditional events, sketched in
\cref{fig:funnel}. The incident must first reach a recording channel, by a
report or by the channel's own detection, and then the police. National Crime
Victimization Survey data indicate that U.S. residents experienced an average
of roughly 250{,}000 hate-crime victimizations a year over 2004--2015, the
majority of which were not reported to the police \citep{masucci2017}. If the
incident reaches the police, an officer must recognise and record an
antisemitic bias motivation, a judgement whose consistency across agencies has
long been questioned \citep{mcdevitt2003}. Finally, the agency must submit the
record to the FBI. Coverage at this last stage is not constant: in the 2021
transition to incident-based reporting, several of the largest agencies were
missing from the initial federal release (\cref{sec:src-police}). Even in an
ordinary year the composition of the reporting panel can move the headline.
From 2022 to 2023 the published national hate-crime total rose from 11{,}634
to 11{,}862 incidents \citep{dojcrs2024}, while on a trend dataset restricted
to agencies with comparable reporting in both years incidents fell by 0.6\%
\citep{fbi2024crimenation}.

Non-governmental audits see a different slice. ADL's Audit includes
non-criminal incidents and draws on reports from victims, law enforcement,
the media and partner organisations \citep{adl2025audit}. Community security
organisations and campus or self-report portals have their own intake
channels, definitions and audiences (\cref{sec:sources}). No single list is
complete, each is selected through a different process, and the lists are not
independent of one another.

\paragraph{Why multiple-systems estimation}
When a closed population is partially enumerated by several overlapping
lists, the pattern of overlaps carries information about the units that every
list missed. This is the logic of capture--recapture, from the two-sample
estimator of \citet{petersen1896} and \citet{lincoln1930} through the
multiple-recapture census of \citet{darroch1958} and the log-linear models of
\citet{fienberg1972}. Under the name \emph{multiple-systems estimation} (MSE)
it has become a standard tool wherever a hidden population leaves traces in
several administrative or documentary sources \citep{bird2018}: killings and
disappearances in armed conflict \citep{lum2013,ball2018,sadinle2018}, victims
of modern slavery \citep{silverman2020} and disease surveillance
\citep{international1995,chao2001}. Antisemitic incidents fit this template:
several organisations separately record overlapping subsets of the same
events, and the dark figure is the unobserved cell of their
cross-classification. One limit is structural. A list can record only
incidents that reach it, so list-based MSE estimates $N_V$, the number of
incidents that reach some recording channel (\cref{sec:funnel}). The
remainder needs a victim survey.

\paragraph{Why antisemitic incidents are hard}
Five features of this domain stress the classical assumptions.
\begin{enumerate}[(i),leftmargin=2.2em]
  \item \emph{Referral dependence.} Lists feed one another. An audit that
  ingests law-enforcement reports is more likely to contain an incident
  because the police recorded it. With two lists this biases the
  Lincoln--Petersen estimator downward (\cref{prop:referral}); with more
  lists its effect depends on how the referral channel sits among them.
  \item \emph{Definitional scope mismatch.} The FBI counts criminal offences
  with a recorded bias motivation. NGO audits also count non-criminal
  harassment, under criteria that change over time. Scope must be harmonised
  before lists are combined.
  \item \emph{Record-linkage error.} The lists share no identifier. Records
  must be matched on sparse, noisy and sensitive attributes (date, place,
  incident type), and false and missed matches distort the capture-history
  table from which the dark figure is estimated \citep{ding1994}.
  \item \emph{Heterogeneous reportability.} An assault at a synagogue and a
  slur on a commuter train do not have the same chance of being recorded
  anywhere. With two lists such heterogeneity leads to underestimation
  \citep{chao1987}; with three or more, even the sign of the bias depends on
  the shape of the visibility distribution (\cref{prop:id-hetero}).
  \item \emph{Coverage shifts.} When the set of reporting agencies changes, as
  in 2021, a change in list coverage can masquerade as a change in incidence.
\end{enumerate}

\paragraph{What this paper does, and does not, do}
This is a methods paper. We develop the statistical framework, prove its
properties, evaluate the estimators in a simulation study on synthetic data
with known ground truth, and pre-register an analysis plan for real sources.
We report \emph{no} empirical estimates of the number of antisemitic incidents
in any jurisdiction or year: we have not linked any real incident data, and
nothing here should be read as a finding about the level or trend of
antisemitism. \Cref{fig:pipeline} summarises the pipeline from raw records to
the dark figure; notation is collected in \cref{tab:notation}.

\IfFileExists{figures/fig-pipeline.tex}{% =====================================================================
%  Figure: the estimation pipeline (LEAD).  Glyphs are schematic.
%  Requires the house colours defined in preamble.tex.
% =====================================================================
\begin{figure}[t]
\centering
\begin{tikzpicture}[x=1cm,y=1cm,
    tile/.style={rounded corners=3pt,draw=LabRule,line width=0.6pt,fill=white,
                 minimum width=2.8cm,minimum height=2.86cm,inner sep=0pt},
    ttl/.style={font=\footnotesize\bfseries,text=LabNavy,align=center,
                anchor=north,inner sep=0pt,text width=2.7cm},
    ref/.style={font=\scriptsize,text=LabGrey,align=center,anchor=south,
                inner sep=0pt,text width=2.7cm},
    flowarrow/.style={-{Stealth[length=4.5pt]},LabNavy,line width=0.8pt},
    glyph/.style={LabNavy,line width=0.6pt},
    lab/.style={font=\scriptsize,text=LabGrey,inner sep=0.5pt},
  ]
  % ------------------------------------------------------------------
  % tile helper:  \tile{name}{x}{y}{title}{reference}
  \newcommand{\tile}[5]{%
    \node[tile] (#1) at (#2,#3) {};
    \node[ttl] at ([yshift=-5pt]#1.north) {#4};
    \node[ref] at ([yshift=4pt]#1.south) {#5};}
  \def\yA{0}\def\yB{-4.05}
  \def\gx{-0.18} % vertical offset of glyph centres within a tile

  % ---------------- row I: from records to links ----------------------
  \tile{t1}{0}{\yA}{$K$ lists of\\raw records}{\S\ref{sec:sources}}
  \tile{t2}{3.3}{\yA}{Within-list\\de-duplication}{\S\ref{sec:linkage}}
  \tile{t3}{6.6}{\yA}{Blocking}{\S\ref{sec:linkage}}
  \tile{t4}{9.9}{\yA}{Fellegi--Sunter /\\Bayesian linkage}{Alg.~\ref{alg:fsem}, Thm.~\ref{thm:review}}
  \tile{t5}{13.2}{\yA}{Posterior over\\linkages $\Lambda$}{\S\ref{sec:linkage}--\ref{sec:models}}

  % ---------------- row II: from links to the dark figure ------------
  \tile{t6}{0}{\yB}{Capture-history\\table $\{n_{\h}\}$}{\S\ref{sec:framework}, Thm.~\ref{thm:linkbias}}
  \tile{t7}{3.3}{\yB}{MSE posterior\\of $N$}{Thm.~\ref{thm:closed}, Eq.~(\ref{eq:lcmcr})}
  \tile{t8}{6.6}{\yB}{Sensitivity to\\the interaction $\omega$}{Cor.~\ref{cor:bounds}}
  \tile{t9}{9.9}{\yB}{Dark figure\\$\nzero=N-n$}{\S\ref{sec:models}, Eq.~(\ref{eq:smallarea})}
  \tile{t10}{13.2}{\yB}{Funnel\\decomposition}{Prop.~\ref{prop:funnel}, Thm.~\ref{thm:shapley}}

  % ---------------- arrows ----------------
  \foreach \a/\b in {t1/t2,t2/t3,t3/t4,t4/t5,t6/t7,t7/t8,t8/t9,t9/t10}{
    \draw[flowarrow] (\a.east) -- (\b.west);}
  \draw[flowarrow,rounded corners=6pt] (t5.south) -- ++(0,-0.62) -| (t6.north);
  \node[fill=white,font=\scriptsize,text=LabNavy,inner sep=3pt] at (6.6,-1.97)
    {draws $\Lambda^{(1)},\dots,\Lambda^{(M)}$\ $\longmapsto$\ tables $n^{(m)}$\ $\longmapsto$\
     $N^{(m)}$:\ linkage uncertainty is propagated, not ignored (Alg.~\ref{alg:pipeline})};

  % ================== glyphs ==================
  % 1. raw records: K = 4 stacks of different sizes
  \begin{scope}[shift={(t1.center)},yshift=\gx cm]
    \foreach \k/\hgt in {1/1.05,2/0.85,3/0.65,4/0.45}{
      \pgfmathsetmacro{\xx}{-0.81+0.43*(\k-1)}
      \draw[glyph,fill=LabTint,rounded corners=1pt] (\xx,-0.5) rectangle ++(0.34,\hgt);
      \foreach \j in {1,...,8}{
        \pgfmathsetmacro{\yy}{-0.5+0.12*\j}
        \pgfmathparse{\yy< -0.5+\hgt-0.05 ? 1 : 0}
        \ifnum\pgfmathresult=1 \draw[LabNavyLight,line width=0.4pt] (\xx+0.06,\yy) -- ++(0.22,0);\fi}
      \node[lab,anchor=north] at (\xx+0.17,-0.52) {$\k$};
    }
  \end{scope}

  % 2. within-list de-duplication: near-duplicates collapse
  \begin{scope}[shift={(t2.center)},yshift=\gx cm]
    \foreach \y/\c in {0.45/LabNavy,0.27/LabNavy,0.0/LabNavyLight,-0.27/LabAccent,-0.45/LabAccent}{
      \fill[\c] (-0.55,\y) circle (2.1pt);}
    \draw[LabGrey,line width=0.4pt,rounded corners=2pt] (-0.7,0.17) rectangle (-0.4,0.56);
    \draw[LabGrey,line width=0.4pt,rounded corners=2pt] (-0.7,-0.56) rectangle (-0.4,-0.17);
    \draw[-{Stealth[length=3pt]},LabGrey,line width=0.5pt] (-0.2,0) -- (0.25,0);
    \foreach \y/\c in {0.3/LabNavy,0.0/LabNavyLight,-0.3/LabAccent}{
      \fill[\c] (0.55,\y) circle (2.1pt);}
  \end{scope}

  % 3. blocking: compare only within blocks
  \begin{scope}[shift={(t3.center)},yshift=\gx cm]
    \def\c{0.17}
    \foreach \i/\j/\n in {0/0/2,2/2/3,5/5/1}{
      \fill[LabNavy!30] (-0.51+\i*\c,0.51-\j*\c) rectangle ++(\n*\c,-\n*\c);}
    \draw[LabRule,line width=0.3pt,step=\c,shift={(-0.51,-0.51)}] (0,0) grid (6*\c,6*\c);
    \draw[LabNavy,line width=0.5pt] (-0.51,-0.51) rectangle (0.51,0.51);
    \node[lab,anchor=north] at (0,-0.53) {pairs compared};
  \end{scope}

  % 4. Fellegi--Sunter weights: u- and m-distributions, review band
  \begin{scope}[shift={(t4.center)},yshift=\gx cm]
    \fill[LabAccent!18] (-0.12,-0.45) rectangle (0.32,0.42);
    \draw[LabGrey,line width=0.6pt,domain=-1.15:0.9,samples=50,smooth]
      plot (\x,{-0.45+0.95*exp(-((\x+0.55)^2)/0.09)});
    \draw[LabNavy,line width=0.8pt,domain=-0.4:1.15,samples=50,smooth]
      plot (\x,{-0.45+0.62*exp(-((\x-0.62)^2)/0.07)});
    \draw[LabNavy,line width=0.4pt] (-1.2,-0.45) -- (1.2,-0.45);
    \node[lab,anchor=north] at (0,-0.47) {weight $w=\log m/u$};
    \node[lab,text=LabAccent] at (0.1,0.22) {review};
  \end{scope}

  % 5. posterior over Lambda: a stack of sampled partitions
  \begin{scope}[shift={(t5.center)},yshift=\gx cm]
    \foreach \o in {0.16,0.08}{
      \draw[LabRule,fill=white,rounded corners=1.5pt,line width=0.5pt]
        (-0.75+\o,-0.5+\o) rectangle (0.6+\o,0.45+\o);}
    \draw[LabNavyLight,fill=white,rounded corners=1.5pt,line width=0.5pt]
      (-0.75,-0.5) rectangle (0.6,0.45);
    \foreach \cx/\ys in {-0.45/{0.25,0,-0.25},-0.07/{0.25,-0.05,-0.3},0.3/{0.15,-0.2}}{
      \foreach \yy in \ys {\fill[LabNavy] (\cx,\yy) circle (1.5pt);}}
    \draw[LabNavy,line width=0.6pt] (-0.45,0.25) -- (-0.07,0.25) -- (0.3,0.15);
    \draw[LabNavy,line width=0.6pt] (-0.45,0) -- (-0.07,-0.05);
    \draw[LabAccent,line width=0.6pt,dash pattern=on 1.5pt off 1pt] (-0.07,-0.3) -- (0.3,-0.2);
    \node[lab,anchor=north] at (-0.075,-0.52) {draws $\Lambda^{(m)}$};
  \end{scope}

  % 6. capture histories: a 3-list Venn diagram; outside = dark figure
  \begin{scope}[shift={(t6.center)},yshift=\gx cm]
    \draw[LabNavy,line width=0.5pt,rounded corners=1.5pt] (-1.15,-0.62) rectangle (1.15,0.6);
    \fill[LabNavy,fill opacity=0.13] (-0.2,0.08) circle (0.36) (0.2,0.08) circle (0.36) (0,-0.24) circle (0.36);
    \draw[LabNavy,line width=0.5pt] (-0.2,0.08) circle (0.36) (0.2,0.08) circle (0.36) (0,-0.24) circle (0.36);
    \node[lab,text=LabNavy] at (-0.66,0.3) {$1$};
    \node[lab,text=LabNavy] at (0.66,0.3) {$2$};
    \node[lab,text=LabNavy] at (0.47,-0.48) {$3$};
    \fill[LabNavy] (0,-0.02) circle (0.9pt);
    \node[lab,text=LabAccent,font=\scriptsize] at (0.8,-0.42) {$\nzero\,?$};
  \end{scope}

  % 7. MSE posterior of N (to the right of the observed n)
  \begin{scope}[shift={(t7.center)},yshift=\gx cm]
    \fill[LabNavy!14,domain=-0.15:0.85,samples=40,smooth]
      plot (\x,{-0.45+0.85*exp(-((\x-0.3)^2)/0.12)}) -- (0.85,-0.45) -- (-0.15,-0.45) -- cycle;
    \draw[LabNavy,line width=0.8pt,domain=-0.55:1.15,samples=60,smooth]
      plot (\x,{-0.45+0.85*exp(-((\x-0.3)^2)/0.12)});
    \draw[LabNavy,line width=0.4pt] (-1.15,-0.45) -- (1.2,-0.45);
    \draw[LabAccent,line width=0.7pt] (-0.75,-0.45) -- (-0.75,0.4);
    \node[lab,text=LabAccent,anchor=south] at (-0.75,0.4) {$n$};
    \node[lab,anchor=north] at (0.3,-0.47) {$N$};
  \end{scope}

  % 8. sensitivity: N-hat as a function of omega, with the bound band
  \begin{scope}[shift={(t8.center)},yshift=\gx cm]
    \fill[LabAccent!14] (-0.35,-0.5) rectangle (0.45,0.6);
    \draw[LabNavy,line width=0.4pt,-{Stealth[length=2.5pt]}] (-1.0,-0.5) -- (1.15,-0.5);
    \draw[LabNavy,line width=0.4pt,-{Stealth[length=2.5pt]}] (-1.0,-0.5) -- (-1.0,0.6);
    \draw[LabNavy,line width=0.9pt,domain=-0.95:0.85,samples=40,smooth]
      plot (\x,{-0.38+0.38*exp(1.05*\x)});
    \draw[LabGrey,line width=0.4pt,dash pattern=on 1.5pt off 1pt] (-1.0,{-0.38+0.38*exp(-0.3675)}) -- (-0.35,{-0.38+0.38*exp(-0.3675)});
    \draw[LabGrey,line width=0.4pt,dash pattern=on 1.5pt off 1pt] (-1.0,{-0.38+0.38*exp(0.4725)}) -- (0.45,{-0.38+0.38*exp(0.4725)});
    \node[lab,anchor=north] at (0.05,-0.52) {$\omega\in[\underline\omega,\overline\omega]$};
    \node[lab,anchor=east] at (-1.0,0.45) {$\widehat N$};
  \end{scope}

  % 9. the dark figure: posterior of n_0
  \begin{scope}[shift={(t9.center)},yshift=\gx cm]
    \fill[LabAccent!22,domain=-0.95:1.15,samples=60,smooth]
      plot (\x,{-0.45+0.82*exp(-((\x+0.25)^2)/0.16)}) -- (1.15,-0.45) -- (-0.95,-0.45) -- cycle;
    \draw[LabAccent,line width=0.9pt,domain=-0.95:1.15,samples=60,smooth]
      plot (\x,{-0.45+0.82*exp(-((\x+0.25)^2)/0.16)});
    \draw[LabNavy,line width=0.4pt] (-1.0,-0.45) -- (1.2,-0.45);
    \draw[LabNavy,line width=0.5pt] (-1.0,-0.45) -- ++(0,0.06);
    \node[lab,anchor=north] at (-1.0,-0.47) {$0$};
    \node[lab,anchor=north] at (0.35,-0.47) {$\nzero$ by domain};
  \end{scope}

  % 10. funnel decomposition: Shapley shares (incidence + four stages)
  \begin{scope}[shift={(t10.center)},yshift=\gx cm]
    \def\yo{-0.08}
    \foreach \x/\hh/\c/\l in {-0.96/-0.3/LabNavyLight/{$\theta$},-0.48/0.42/LabNavy/{vic},
                              0/0.28/LabNavyMid/{pol},0.48/0.6/LabAccent/{cls},
                              0.96/0.2/LabNavyLight/{sub}}{
      \fill[\c] (\x-0.11,\yo) rectangle ++(0.22,\hh);
      \node[lab,anchor=north] at (\x,-0.42) {\l};}
    \draw[LabNavy,line width=0.4pt] (-1.15,\yo) -- (1.15,\yo);
  \end{scope}
\end{tikzpicture}
\caption{\textbf{From records to the dark figure.} Top row: records from
$K$ lists are de-duplicated within lists, restricted by blocking to plausible
candidate pairs, and linked across lists with a Fellegi--Sunter or Bayesian
model; uncertain pairs (orange band) can be sent to clerical review. The
output is a \emph{posterior distribution} over linkage structures $\Lambda$
rather than a single matching. Bottom row: each draw $\Lambda^{(m)}$ induces a
capture-history table; multiple-systems estimation yields a posterior for
$N$; the untestable $K$-way interaction $\omega$ is varied over a stated range;
and the dark figure $\nzero=N-n$ (the region outside every list) is reported
by domain and decomposed along the reporting funnel of \cref{fig:funnel} into
Shapley shares of incidence ($\theta$) and the four stages ($\rho_{\mathrm{vic}},\dots,\rho_{\mathrm{sub}}$).
Labels under each tile point to the corresponding results.}
\label{fig:pipeline}
\end{figure}
}{}%
\IfFileExists{tables/tab-notation.tex}{% =====================================================================
%  Table: notation (LEAD).  Symbols follow SPEC.md exactly.
% =====================================================================
\begin{table}[tbp]
\caption{\textbf{Notation.} Symbols used throughout. Bold $\h$ denotes a
capture history; $\zero$ is the all-zero history.}
\label{tab:notation}
\small
\renewcommand{\arraystretch}{1.08}
\newcommand{\nhead}[1]{\multicolumn{2}{@{}l}{\footnotesize\scshape\bfseries\color{LabNavy}#1}\\[1pt]}
\begin{minipage}[t]{0.485\textwidth}
\begin{tabularx}{\linewidth}[t]{@{}>{$}l<{$}@{\hspace{0.8em}}>{\raggedright\arraybackslash}X@{}}
\toprule
\nhead{Population and lists}
d & domain (state $s$ $\times$ year $t$) \\
U_d,\ N & in-scope antisemitic incidents that reach a recording channel; unknown size (\cref{sec:funnel} writes $U_d^V$, $N_V$, and uses $U_d$, $N$ for all incidents) \\
k\in\Kset & list; $K=4$: police (FBI), NGO audit, community security org., campus/self-report portal \\
Y_{ik} & $1$ if incident $i$ is on list $k$ \\
\midrule
\nhead{Capture histories and counts}
\h\in\Hset=\{0,1\}^K & capture history; $|\h|=\sum_k h_k$ \\
n_{\h} & incidents with history $\h$ \\
n=\sum_{\h\neq\zero}n_{\h} & distinct incidents observed \\
\nzero=N-n & \textbf{the dark figure} \\
n_{k+},\ f_j & size of list $k$; incidents on exactly $j$ lists \\
D_k=N/n_{k+} & dark-figure ratio relative to list $k$ \\
s_{\h}=(-1)^{|\h|+1} & alternating sign of a history \\
\mathcal A(\bm x)=\prod_{\h\neq\zero}x_{\h}^{s_{\h}} & alternating product \\
\midrule
\nhead{Reporting funnel}
\rvic & reaches a recording channel, by a report or by the channel's own detection \\
\rpol & reaches the police $\mid$ reaches a channel \\
\rcls & classified as antisemitic $\mid$ reaches the police \\
\rsub & submitted to the FBI $\mid$ classified \\
N_V & incidents that reach a channel; what list-based MSE estimates \\
\nFBI & FBI count; $\E\,\nFBI=\E[N_V]\,\rpol\rcls\rsub$ \\
\bottomrule
\end{tabularx}
\end{minipage}\hfill
\begin{minipage}[t]{0.485\textwidth}
\begin{tabularx}{\linewidth}[t]{@{}>{$}l<{$}@{\hspace{0.8em}}>{\raggedright\arraybackslash}X@{}}
\toprule
\nhead{Capture models}
\mu_{\h}=\E[n_{\h}] & expected cell count \\
u_S,\ S\subseteq\Kset & log-linear terms: $\log\mu_{\h}=\sum_{S\subseteq\supp(\h)}u_S$ \\
\omega=u_{\Kset} & highest-order interaction; sensitivity parameter \\
\bar\omega=(-1)^K\omega & oriented interaction: $\mu_{\zero}=e^{\bar\omega}\mathcal A(\bm\mu_+)$ \\
p_{ik} & $\Pr(Y_{ik}=1\mid x_i)$, capture probability \\
z_i,\ \pi_j,\ \lambda_{jk} & latent class; class weight; class-$j$ capture probability on list $k$ \\
\midrule
\nhead{Record linkage}
r\in R_k,\ a_r & record of list $k$; its attributes \\
\Link,\ \Link^{\star},\ \Linkhat & linkage structure (partition of records into incidents); truth; estimate \\
\gamma_{rr'} & comparison vector of a record pair \\
m(\gamma),\ u(\gamma) & $\Pr(\gamma\mid\text{match})$, $\Pr(\gamma\mid\text{non-match})$ \\
w=\log(m/u) & Fellegi--Sunter weight \\
\delta,\ \varepsilon & per-record probabilities of detachment and of a false merge under LE (pairwise false non-match rate $1-(1-\delta)^2$) \\
\Lop(\delta,\varepsilon) & linkage operator on capture-history tables \\
\midrule
\nhead{Small-area model}
\theta_{st} & incidence per unit exposure, state $s$, year $t$ \\
b_s,\ \tau_t,\ \chi_{st} & state effect; national trend; state-specific trend \\
c_{st,k} & coverage of list $k$ (police: agency participation) \\
\bottomrule
\end{tabularx}
\end{minipage}
\end{table}
}{}%
\paragraph{Contributions}
Our contributions are a synthesis, generalisation and application of
existing ideas to a new and sensitive domain. Classical results are
attributed where they are used.
\begin{enumerate}[label=\textbf{C\arabic*},ref=C\arabic*,leftmargin=2.4em,itemsep=3pt]
  \item\label{contrib:id} \textbf{Identification and sensitivity.} Without a
  restriction on the highest-order interaction $\omega$ of the log-linear
  model, the dark figure is not identified (\cref{thm:nonid}; compare
  \citealp{aleshinguendel2024}). Without a $K$-way interaction, the
  maximum-likelihood estimate of the dark figure is an alternating product of
  cell counts over the Boolean lattice $\{0,1\}^K$
  \citep{fienberg1972,bishop1975}. We retain $\omega$ in this classical closed
  form as an explicit sensitivity parameter, which enters as the factor
  $e^{\pm\omega}$ (\cref{thm:closed}, \cref{fig:lattice}) and makes
  offset-based sensitivity analysis \citep{gerritse2015} analytic. The bounds
  are sharp (\cref{cor:bounds}), and the direction of bias is the sign of an
  oriented interaction (\cref{prop:direction}). With three or more lists,
  neither positive dependence nor heterogeneous visibility fixes that sign
  (\cref{prop:id-hetero}); in a stylised three-list example the model without
  a $K$-way interaction overstates the dark figure more than twelvefold
  (\cref{ex:id-hetero}). Sensitivity analysis must therefore be two-sided
  unless a known mechanism fixes the sign. Chao's lower bound
  \citep{chao1987} holds under heterogeneity across incidents, but
  heterogeneity across lists reverses it (\cref{thm:chao,prop:id-mt}), so we
  use it as a diagnostic. We model referral dependence and show how provenance
  metadata restore identification (\cref{prop:referral}). A sensitivity
  frontier states how much unmeasured dependence a claim such as ``the dark
  figure exceeds the police count'' can survive (\cref{thm:id-frontier}).
  \item\label{contrib:link} \textbf{Linkage error as a linear operator.}
  Following the dual-system analysis of matching error by \citet{ding1994}
  and its multi-list extensions \citep{ding1996,diconsiglio2018}, we
  represent imperfect record linkage under a stylised error model (LE) as an
  operator $\Lop(\delta,\varepsilon)$ on the true capture-history table, where
  $\delta$ and $\varepsilon$ are per-record probabilities of a missed and a
  false match. We derive first-order bias coefficients for the $K$-list
  estimator (\cref{thm:linkbias}). Under independent capture, false matches do
  more damage when lists are sparse and missed matches when they are dense,
  with a closed-form crossover (\cref{cor:link-indep,rem:link-reading}).
  Missed matches inflate the two-list dark figure, but with three or more
  strongly dependent lists that direction can reverse
  (\cref{rem:link-reversal}). We derive the correction implied under LE
  (\cref{cor:linkcorrect}) and a total-variance decomposition that quantifies
  how far plug-in intervals undercover (\cref{prop:totalvar}). LE has no move
  for a swap, a record linked to the wrong partner, and our simulation shows
  that the correction over-corrects when swaps dominate. Mixing over linkage
  draws is therefore our primary approach, and the correction is a
  diagnostic. We also design clerical review and a random audit of linkage
  quality (\cref{thm:review,alg:review,prop:link-audit}).
  \item\label{contrib:bayes} \textbf{A Bayesian model.} We assemble the
  latent-class MSE of \citet{manriquevallier2016}, with a Dirichlet-process
  prior (\cref{eq:lcmcr}), covariate-dependent capture, and a spatiotemporal
  small-area model for state--year counts (\cref{eq:smallarea}) in which
  police-list coverage enters as a known offset computed from agency
  participation, so that breaks such as the 2021 transition are absorbed by
  data. The small-area model is specified, not evaluated. Linkage uncertainty
  is propagated through a modular (``cut'') posterior that mixes over draws of
  the linkage structure (\cref{eq:models-cut}, \cref{alg:pipeline}), in the
  spirit of \citet{tancredi2011} and \citet{sadinle2018}.
  \item\label{contrib:funnel} \textbf{Where antisemitic incidents go missing.}
  We write the police count as a multiplicative reporting funnel
  (\cref{prop:funnel}, \cref{fig:funnel}) and state which data identify each
  stage. No list-based method can see incidents that never reach a recording
  channel, so that stage requires a victim survey. A Shapley-value
  decomposition attributes the dispersion of recorded rates across
  jurisdictions, or the gap between two of them, to incidence and to each
  reporting stage (\cref{thm:shapley}).
  \item\label{contrib:sim} \textbf{A simulation study.} In synthetic worlds
  with heterogeneous reportability, referral dependence and noisy records, we
  compare Lincoln--Petersen, log-linear model averaging, latent-class DP-MSE
  and Chao's estimator, and naive, linkage-corrected and
  uncertainty-propagating pipelines, on bias, root-mean-squared error and
  interval coverage (\cref{sec:simulation}). Simulated records carry
  realistic attribute noise and are linked by FS-EM (\cref{alg:fsem}); LE
  enters only through the correction, which is thereby tested under
  misspecification.
  \item\label{contrib:plan} \textbf{A pre-registered analysis plan.} We fix
  in advance how FBI hate-crime data, the ADL Audit, community security
  organisations, campus data and the NCVS would be harmonised, linked under
  privacy constraints and analysed, with New York City as the candidate pilot
  (\cref{sec:plan}), together with an ethics protocol (\cref{sec:ethics}).
\end{enumerate}

\paragraph{Related work}
The effect of matching error on dual-system estimates was analysed by
\citet{ding1994}, with multi-list extensions by \citet{ding1996} and
\citet{diconsiglio2018}. Joint Bayesian models of record linkage and
population size were developed by \citet{tancredi2011}. \citet{sadinle2018}
showed how to propagate linkage uncertainty into MSE for human-rights data,
building on Bayesian bipartite matching \citep{sadinle2017} and the
record-linkage theory of \citet{fellegi1969}. Heterogeneity in capture
probabilities has been addressed through latent-class and Dirichlet-process
models \citep{manriquevallier2016} and through lower bounds
\citep{chao1987}. On hate crime, the under-reporting literature documents the
gap between victimisation and police statistics \citep{pezzella2019} and the
fragility of bias classification \citep{mcdevitt2003}. As far as we know,
these strands have not been brought together for antisemitic incidents, where
referral between lists, definitional mismatch and shifting police coverage
occur together.

\paragraph{Roadmap}
\Cref{sec:sources} describes the candidate lists and \cref{sec:framework} the
population, lists and capture model. \Cref{sec:identification} develops
identification, sensitivity bounds and referral dependence, \cref{sec:linkage}
record-linkage error and clerical review, and \cref{sec:models} the Bayesian
model. \Cref{sec:funnel} develops the reporting funnel, \cref{sec:simulation}
reports the simulation study, \cref{sec:plan} gives the pre-registered
analysis plan and \cref{sec:ethics} the ethics protocol. \Cref{sec:discussion}
discusses limitations. Proofs are in \cref{app:idproofs,app:linkproofs}, with
the remaining funnel proof in \cref{app:bayes-proofs}; computational and
simulation details are in \cref{app:computation,app:simdetails}.
